\documentclass[10pt]{article}
\usepackage[T1]{fontenc}
\usepackage[utf8]{inputenc}
\usepackage[french]{babel}
\usepackage{amsmath,amssymb,mathrsfs,mathpazo}
\usepackage{geometry}
\geometry{a4paper,margin=2cm}
\usepackage{graphicx,array,tabularx,tabularray,colortbl,multirow,multicol,booktabs,enumitem,pifont,ulem,fancybox,diagbox,cancel}
\usepackage{pstricks,pstricks-add,pst-plot,pst-tree,pst-node,pst-text,pst-func,pst-math,pst-bspline,pst-3dplot,pst-eucl,pst-solides3d}
\let\pscalculate\undefined
\usepackage{pst-bezier}
\usepackage{tikz,pgfplots,tkz-tab}
\usetikzlibrary{arrows,arrows.meta,patterns,calc,positioning,trees,decorations.markings,shapes.geometric,shapes.misc,shapes.arrows,intersections,decorations.pathmorphing,decorations.pathreplacing,backgrounds,shadows,babel}
\pgfplotsset{compat=1.18}
\usepgfplotslibrary{fillbetween}
\usepackage[np]{numprint}
\usepackage{listings,makecell,pas-tableur,twoopt}
\let\barre\undefined
\usepackage[pstricks]{tablvar}
\definecolor{darkgreen}{RGB}{0,100,0}
\providecommand{\red}{\color{red}}
\providecommand{\blue}{\color{blue}}
\providecommand{\green}{\color{green}}
\providecommand{\black}{\color{black}}
\providecommand{\white}{\color{white}}
\providecommand{\gray}{\color{gray}}
\providecommand{\cyan}{\color{cyan}}
\providecommand{\magenta}{\color{magenta}}
\providecommand{\yellow}{\color{yellow}}
\providecommand{\violet}{\color{violet}}
\providecommand{\orange}{\color{orange}}
\providecommand{\euro}{\texteuro}
\providecommand{\EUR}{\texteuro}
\providecommand{\ieme}{\textsuperscript{e}}
\providecommand{\degres}{\ensuremath{^\circ}}
\providecommand{\degree}{\ensuremath{^\circ}}
\providecommand{\vv}[1]{\overrightarrow{#1}}
\providecommand{\Coord}[1]{#1}
\providecommand{\decosix}{}
\providecommand{\starredbullet}{\ensuremath{\bullet}}
\pagestyle{empty}
\setlength{\parindent}{0pt}
\usepackage{adjustbox}\DeclareUnicodeCharacter{2236}{:}\DeclareUnicodeCharacter{2190}{\ensuremath{\leftarrow}}
\providecommand{\R}{\mathbb{R}}
\providecommand{\N}{\mathbb{N}}
\providecommand{\Z}{\mathbb{Z}}
\providecommand{\Q}{\mathbb{Q}}
\providecommand{\D}{\mathbb{D}}
\providecommand{\C}{\mathbb{C}}
\providecommand{\e}{{\mathrm{e}}}
\providecommand{\vect}[1]{\overrightarrow{#1}}
\providecommand{\vectt}[1]{\overrightarrow{\mathrm{#1}}}
\providecommand{\vvt}[1]{\overrightarrow{\mathrm{#1}}}
\providecommand{\np}[2][]{{#2~#1}}
\providecommand{\vv}[1]{{\overrightarrow{#1}}}
\renewcommand{\d}{{\mathrm{d}}}
\providecommand{\degres}{{^\circ}}
\providecommand{\degre}{{^\circ}}
\providecommand{\blue}{\color{blue}}
\providecommand{\red}{\color{red}}
\providecommand{\enskip}{\quad}
\providecommand{\llbracket}{\lbrack\!\lbrack}
\providecommand{\rrbracket}{\rbrack\!\rbrack}
\providecommand{\euro}{{\text{€}}}
\providecommand{\no}{{\text{n°}}}
\providecommand{\interval}[2]{{[#1;#2]}}
\providecommand{\intervFO}[2]{{[#1;#2[}}
\providecommand{\dotfill}{\ldots}

\providecommand{\ds}{\displaystyle}
\providecommand{\textding}[1]{#1}
\providecommand{\Oij}{\(\mathrm{O};\vec{\imath},\vec{\jmath}\)}
\providecommand{\Oijk}{\(\mathrm{O};\vec{\imath},\vec{\jmath},\vec{k}\)}
\providecommand{\Ouv}{\(\mathrm{O};\vec{u},\vec{v}\)}
\providecommand{\intervFF}[2]{\ensuremath{[#1;#2]}}
\renewcommand{\euro}{\ifmmode\text{\texteuro}\else\texteuro\fi}
\DeclareUnicodeCharacter{00B0}{\ifmmode^{\circ}\else\textdegree\fi}
\DeclareUnicodeCharacter{20AC}{\ifmmode\text{\texteuro}\else\texteuro\fi}
% Énoncés : Éducation nationale. Transcription et corrigés : APMEP, Denis Vergès et contributeurs. Découpage et classement : CoopMaths, Sébastien Lozano, Éric Elter et contributeurs. Adaptation de la présentation : Maths Explorer.
% CC BY-SA : https://coopmaths.fr/www/about/
% https://coopmaths.fr/alea/static/eam/2026/tex/eam_specifique_2026_04_polynesie_1_cor.tex
\begin{document}
\noindent\textit{APMEP / CoopMaths. CC BY-SA. Présentation adaptée par Maths Explorer.}\par\medskip
\begin{enumerate}
	\item
	\begin{enumerate}
		\item Diminuer de 2\,\% revient à multiplier par $1 - \dfrac{2}{100} = 0,98$.
		
		$u_1 = 0,98 \times \np{60000} = \np{58800}$.

		On prévoit \np{58800} naissances en 2026.

		\item $v_1 = \np{40000} + 900 = \np{40900}$.

		On prévoit \np{40900} décès en 2026.
	\end{enumerate}

	\item
	\begin{enumerate}
		\item D'une année à la suivante, le nombre de naissances diminue de 2\,\% donc, pour tout entier naturel $n$ :
		\[u_{n+1} = u_n - \dfrac{2}{100}u_n = u_n(1 - 0,02) = 0,98u_n.\]

		\item On passe d'un terme au suivant en multipliant toujours par le même nombre $0,98$ : la suite
		$(u_n)$ est \textbf{géométrique} de raison $q = 0,98$ et de premier terme $u_0 = \np{60000}$.

		\item Pour tout entier naturel $n$, $u_n = u_0 \times q^n$.
		
		 Pour tout entier naturel $n$, $u_n = \np{60000} \times 0,98^n$.
	\end{enumerate}

	\item
	\begin{enumerate}
		\item Pour tout entier naturel $n$, $v_{n+1} = v_n + 900$. La suite $(v_n)$
		est \textbf{arithmétique} de raison $r = 900$ et de premier terme $v_0 = \np{40000}$.

		\item Pour tout entier naturel $n$,  $v_n = v_0 + nr$
		
		Pour tout entier naturel $n$, $v_n = \np{40000} + 900n$.
	\end{enumerate}

	\item
	\begin{enumerate}
		\item On saisit en C3 la formule \texttt{=C2+900}.

		\item La ligne 12 correspond à $n = 10$ :
		$v_{10} = \np{40000} + 900 \times 10 = \np{49000}$.

		La cellule C12 affiche \np{49000}.
	\end{enumerate}

	\item Le solde naturel de l'année $2025+n$ est $u_n - v_n$ ; il est négatif lorsque $u_n < v_n$.

		Dans la tableau donné, on remarque que pour $n = 10$, $u_{10} > v_{10}$ et pour $n=11$, $u_{11} < v_{11}$.
	
	Comme $(u_n)$ est décroissante et $(v_n)$ croissante, le
	solde reste ensuite négatif.

	Le solde naturel devient négatif à partir de l'année
	$2025 + 11 = 2036$.
\end{enumerate}

\end{document}
